En esta práctica se analizó la relación entre la presión hidrostática y la profundidad del agua para determinar experimentalmente el valor de la gravedad. Utilizando un manómetro conectado a una manguera, se registraron los cambios de presión a medida que esta se sumergía en el tanque. Al graficar los datos y hacer una regresión lineal, la pendiente obtenida permitió calcular una gravedad de $9.4 \pm 0.1\ \si{m/s^2}$. Este resultado tiene un error relativo del $-4\%$ comparado con la gravedad esperada en Cali. Las diferencias se explican por la dificultad para establecer el cero en el manómetro y por el volumen extra que desplazaban los instrumentos al entrar al agua.
